\chapter{Laurent Lafforgue --- Fields Medal 2002}\label{ch2002laurentlafforgue:laurent-lafforgue-fields-medal-2002-a-comprehensive-report-on-his-work-and-impact}

\begin{center}\rule{0.5\linewidth}{0.5pt}\end{center}

\textbf{Prepared:} July 2026 \textbf{Subject:} Laurent Lafforgue (b.
November 6, 1966, Antony, France) \textbf{Position:} Professor, Institut
des Hautes Études Scientifiques (IHÉS), Bures-sur-Yvette (since 2000);
CNRS researcher (2001--) \textbf{Doctoral Advisor:} Gérard Laumon
(Université Paris-Sud) \textbf{Fields Medal Citation:} \emph{``For his
proof of the Langlands correspondence for GL(r) over function fields.''}

\begin{center}\rule{0.5\linewidth}{0.5pt}\end{center}

\section{1. Introduction}\label{ch2002laurentlafforgue:introduction}

Laurent Lafforgue was awarded the Fields Medal in 2002 for proving the
\textbf{Langlands correspondence for GL(r) over function fields}. This
was a significant achievement in number theory and
representation theory of the late twentieth century.

\begin{prereqs} Algebraic geometry (stacks, moduli spaces), number theory (automorphic forms, L-functions), basic representation theory. \end{prereqs}

\begin{objectives} Explain Lafforgue's proof of the Langlands correspondence for function fields; describe the moduli space of shtukas; understand the role of étale cohomology. \end{objectives}


The Langlands program, initiated by Robert Langlands in the 1960s, is a
sweeping web of conjectures relating number theory (Galois groups) and
representation theory (automorphic forms). Proving these conjectures has
been one of the main drivers of research in number theory for over fifty
years. Lafforgue's proof for GL(r) over function fields was a monumental
achievement that required the development of new geometric and
cohomological techniques.

This report provides a comprehensive account of Lafforgue's work, its
mathematical context, and its significant impact on arithmetic geometry and
the Langlands program.

\begin{center}\rule{0.5\linewidth}{0.5pt}\end{center}

\section{2. Biographical Background}\label{ch2002laurentlafforgue:biographical-background}

\subsection{2.1 Early Life}\label{ch2002laurentlafforgue:early-life}

Laurent Lafforgue was born in Antony, France in 1966. He showed early
mathematical talent and entered the Lycée Louis-le-Grand in Paris for
his preparatory classes. He was admitted to the École Polytechnique in
1986 and later to the École Normale Supérieure (ENS) in 1988.

\subsection{2.2 Education}\label{ch2002laurentlafforgue:education}

At ENS, Lafforgue studied under Gérard Laumon, a leading expert in
automorphic forms and the Langlands program. His PhD thesis, completed
in 1994 at the Université Paris-Sud, was titled ``Chtoucas de Drinfeld
et conjecture de Ramanujan--Petersson'' (Drinfeld shtukas and the
Ramanujan--Petersson conjecture).

\subsection{2.3 Academic Career}\label{ch2002laurentlafforgue:academic-career}

\begin{itemize}
\tightlist
\item
 \textbf{1994--2000:} CNRS researcher, Université Paris-Sud.
\item
 \textbf{2000--present:} Professor, Institut des Hautes Études
 Scientifiques (IHÉS).
\item
 \textbf{2001--2017:} Also CNRS researcher.
\item
 \textbf{2002:} Fields Medal, International Congress of Mathematicians,
 Beijing.
\end{itemize}

\begin{center}\rule{0.5\linewidth}{0.5pt}\end{center}

\section{3. The Langlands Program}\label{ch2002laurentlafforgue:the-langlands-program}

\subsection{3.1 The Vision}\label{ch2002laurentlafforgue:the-vision}

The Langlands program, proposed by Robert Langlands in a famous 1967
letter to André Weil, consists of a series of conjectures connecting: -
\textbf{Galois groups:} The absolute Galois group of a number field or
function field. - \textbf{Automorphic forms:} Harmonic analysis on
adelic groups. - \textbf{L-functions:} Generalizations of the Riemann
zeta function.

The central conjecture (the \textbf{Langlands correspondence}) states
that there is a bijection between: - Irreducible \(n\)-dimensional
representations of the Galois group. - Cuspidal automorphic
representations of GL(\(n\)) over the adeles.

\subsection{3.2 The Role of L-Functions}\label{ch2002laurentlafforgue:the-role-of-l-functions}

The correspondence is expected to preserve L-functions: the Artin
L-function of a Galois representation should equal the automorphic
L-function of the corresponding automorphic representation. This would
imply the Artin conjecture (that Artin L-functions are entire).

\subsection{3.3 Function Fields vs.~Number
Fields}\label{ch2002laurentlafforgue:function-fields-vs.-number-fields}

The Langlands program can be formulated for both number fields (finite
extensions of \(\mathbb{Q}\)) and function fields (finite extensions of
\(\mathbb{F}_q(T)\)). The function field case is more tractable because:
- \textbf{Geometric methods:} Function fields are fields of rational
functions on algebraic curves over finite fields. - \textbf{The Weil
conjectures:} Proved by Deligne, providing powerful cohomological tools.
- \textbf{Drinfeld's shtukas:} A geometric construction of automorphic
forms.

\begin{center}\rule{0.5\linewidth}{0.5pt}\end{center}

\section{4. Drinfeld's Work}\label{ch2002laurentlafforgue:drinfelds-work}

\subsection{4.1 Drinfeld Shtukas}\label{ch2002laurentlafforgue:drinfeld-shtukas}

Vladimir Drinfeld introduced the concept of \textbf{shtukas} (a Russian
word meaning ``thing'') in the 1970s. A shtuka is a vector bundle with
additional structure --- essentially, a family of vector bundles on a
curve equipped with a Frobenius morphism.

Drinfeld used shtukas to prove the Langlands correspondence for GL(2)
over function fields, for which he received the Fields Medal in 1990. He
also introduced the \textbf{elliptic sheaves} and \textbf{Drinfeld
modules} that generalize the theory of elliptic curves to function
fields.

\subsection{4.2 Beyond GL(2)}\label{ch2002laurentlafforgue:beyond-gl2}

Extending Drinfeld's methods from GL(2) to GL(\(r\)) for arbitrary \(r\)
required: - A theory of \textbf{moduli spaces of shtukas} for higher
rank. - A \textbf{trace formula} for computing the cohomology of these
moduli spaces. - A proof that the cuspidal automorphic representations
appear in the cohomology.

This is precisely what Lafforgue accomplished.

\begin{center}\rule{0.5\linewidth}{0.5pt}\end{center}

\section{5. The Geometrization of the Langlands
Program}\label{ch2002laurentlafforgue:the-geometrization-of-the-langlands-program}

\subsection{5.1 The Geometric Langlands
Program}\label{ch2002laurentlafforgue:the-geometric-langlands-program}

The \textbf{geometric Langlands program}, initiated by Beilinson and
Drinfeld, is a geometric reformulation of the Langlands correspondence.
Instead of automorphic forms on adelic groups, one considers perverse
sheaves on moduli spaces of bundles.

Lafforgue's work is part of this geometric approach. He uses the moduli
space of shtukas as a geometric object whose cohomology realizes the
Langlands correspondence.

\subsection{5.2 Chtoucas}\label{ch2002laurentlafforgue:chtoucas}

Lafforgue uses the French term \textbf{chtoucas} (a transliteration of
the Russian ``shtuka'') to emphasize the geometric nature of the
construction. A chtouca of rank \(r\) on a curve \(X\) over
\(\mathbb{F}_q\) consists of: - A vector bundle \(E\) of rank \(r\) on
\(X\). - A Frobenius morphism \(\text{Fr}: E \to E\). - A modification
of \(E\) at a finite set of points.

\begin{center}\rule{0.5\linewidth}{0.5pt}\end{center}

\section{6. The Moduli Space of
Shtukas}\label{ch2002laurentlafforgue:the-moduli-space-of-shtukas}

\subsection{6.1 Definition}\label{ch2002laurentlafforgue:definition}

For a smooth projective curve \(X\) over \(\mathbb{F}_q\), the moduli
space \(\text{Cht}^r_{k,m}\) parametrizes shtukas of rank \(r\) with
\(k\) modifications (or ``legs''). The legs are points of \(X\) where
the Frobenius morphism is allowed to have poles.

\subsection{6.2 Properties}\label{ch2002laurentlafforgue:properties}

The moduli space of shtukas is: - \textbf{A Deligne--Mumford stack} (an
algebraic stack with finite stabilizers). - \textbf{Proper} over \(X^k\)
(the product of \(k\) copies of \(X\)). - \textbf{Smooth} of dimension
\(2r(k-1)\). - \textbf{Equipped with a Hecke action} for each point of
\(X\).

\subsection{6.3 The Hecke Action}\label{ch2002laurentlafforgue:the-hecke-action}

At each point \(x \in X\), there is a Hecke correspondence: - A
correspondence on the moduli space that adds a modification at \(x\). -
This action generates an algebra of Hecke operators. - The eigenvalues
of these operators correspond to automorphic representations.

\begin{center}\rule{0.5\linewidth}{0.5pt}\end{center}

\section{7. The Cohomology of the Moduli
Space}\label{ch2002laurentlafforgue:the-cohomology-of-the-moduli-space}

\subsection{7.1 The Trace Formula}\label{ch2002laurentlafforgue:the-trace-formula}

Lafforgue uses the \textbf{Grothendieck--Lefschetz trace formula} to
compute the cohomology of the moduli space of shtukas. The trace formula
expresses the number of fixed points of a Frobenius morphism as an
alternating sum of traces on cohomology.

For the moduli space of shtukas, the number of points over
\(\mathbb{F}_{q^n}\) can be computed in two ways: 1.
\textbf{Geometrically:} By counting shtukas over the finite field. 2.
\textbf{Cohomologically:} By computing the trace of Frobenius on the
cohomology.

\subsection{7.2 The Main Term}\label{ch2002laurentlafforgue:the-main-term}

The main term in the trace formula gives:

\[
\sum_{\pi} \prod_{i=1}^k \text{tr}(\pi(\text{Fr}_{x_i}))
\]

where the sum is over cuspidal automorphic representations \(\pi\) of
GL(\(r\)) over the function field, and \(\pi(\text{Fr}_{x_i})\) are the
Hecke operators at the points \(x_i\).

\subsection{7.3 The Geometric Side}\label{ch2002laurentlafforgue:the-geometric-side}

The geometric side counts shtukas over \(\mathbb{F}_{q^n}\). By
analyzing the fixed points of Frobenius, Lafforgue shows that the
geometric side matches the automorphic side, establishing the Langlands
correspondence.

\begin{center}\rule{0.5\linewidth}{0.5pt}\end{center}

\section{8. The Compactification of the Moduli
Space}\label{ch2002laurentlafforgue:the-compactification-of-the-moduli-space}

\subsection{8.1 The Need for
Compactification}\label{ch2002laurentlafforgue:the-need-for-compactification}

The moduli space of shtukas is not compact. To apply the trace formula
in cohomology, one needs a compactification that: - Is proper
(complete). - Has a well-behaved cohomology. - Admits an action of the
Hecke operators.

\subsection{8.2 Lafforgue's
Compactification}\label{ch2002laurentlafforgue:lafforgues-compactification}

Lafforgue constructed a compactification of the moduli space of shtukas
using: - \textbf{Toric geometry:} Gluing in boundary divisors
corresponding to degenerate shtukas. - \textbf{The Hilbert--Mumford
criterion:} A stability condition that ensures the compactified space is
proper. - \textbf{The theory of algebraic stacks:} To handle the moduli
problem properly.

The compactification is a smooth proper algebraic stack that contains
the original moduli space as an open dense substack.

\begin{center}\rule{0.5\linewidth}{0.5pt}\end{center}

\section{9. The Stabilization of the Trace
Formula}\label{ch2002laurentlafforgue:the-stabilization-of-the-trace-formula}

\subsection{9.1 Arthur's Trace Formula}\label{ch2002laurentlafforgue:arthurs-trace-formula}

The Arthur--Selberg trace formula is an effective tool for studying
automorphic forms. For GL(\(r\)) over a function field, the trace
formula relates: - \textbf{The spectral side:} Sums over automorphic
representations. - \textbf{The geometric side:} Sums over conjugacy
classes in the group.

\subsection{9.2 Lafforgue's Contribution}\label{ch2002laurentlafforgue:lafforgues-contribution}

Lafforgue proved a \textbf{stabilization} of the trace formula for
GL(\(r\)) over function fields. This means that he expressed the trace
formula as a sum over ``stable'' conjugacy classes, which are easier to
analyze geometrically.

The stabilization is essential for the proof of the Langlands
correspondence because it allows the geometric side to be matched with
the spectral side.

\begin{center}\rule{0.5\linewidth}{0.5pt}\end{center}

\section{10. The Ramanujan--Petersson
Conjecture}\label{ch2002laurentlafforgue:the-ramanujanpetersson-conjecture}

\subsection{10.1 The Conjecture}\label{ch2002laurentlafforgue:the-conjecture}

The Ramanujan--Petersson conjecture bounds the size of the Hecke
eigenvalues of cuspidal automorphic forms. For GL(\(r\)) over function
fields, it states that:

\[
|\alpha_{i,\pi}(x)| = q^{1/2} \quad \text{for all } i, x
\]

where \(\alpha_{i,\pi}(x)\) are the Satake parameters of \(\pi\) at a
place \(x\).

\subsection{10.2 Lafforgue's Proof}\label{ch2002laurentlafforgue:lafforgues-proof}

Lafforgue proved the Ramanujan--Petersson conjecture for GL(\(r\)) over
function fields as a consequence of his proof of the Langlands
correspondence. The bound follows from the fact that the eigenvalues of
Frobenius on the cohomology of the moduli space of shtukas are Weil
numbers (by Deligne's theorem on the Weil conjectures).

\begin{center}\rule{0.5\linewidth}{0.5pt}\end{center}

\section{11. The Global Langlands
Correspondence}\label{ch2002laurentlafforgue:the-global-langlands-correspondence}

\subsection{11.1 The Main Theorem}\label{ch2002laurentlafforgue:the-main-theorem}

Lafforgue proved:

\begin{quote}
\textbf{Theorem (Lafforgue, 2002):} For a function field \(F\) over a
finite field \(\mathbb{F}_q\), there is a bijection between: -
Irreducible \(r\)-dimensional representations of the absolute Galois
group of \(F\). - Cuspidal automorphic representations of GL(\(r\)) over
\(F\).

The bijection preserves L-functions and \(\epsilon\)-factors.
\end{quote}

\subsection{11.2 The Bijection}\label{ch2002laurentlafforgue:the-bijection}

The bijection is given by:

\[
\sigma \leftrightarrow \pi(\sigma)
\]

where: - \(\sigma\) is an \(r\)-dimensional \(\ell\)-adic representation
of \(\text{Gal}(\overline{F}/F)\). - \(\pi(\sigma)\) is the
corresponding cuspidal automorphic representation of
GL(\(r, \mathbb{A}_F\)).

\subsection{11.3 Compatibility}\label{ch2002laurentlafforgue:compatibility}

The correspondence is compatible with: - \textbf{L-functions:}
\(L(s, \sigma) = L(s, \pi(\sigma))\). - \textbf{Local Langlands:} At
each place of \(F\), the restriction of \(\sigma\) to the decomposition
group corresponds to the local component of \(\pi(\sigma)\). -
\textbf{Functoriality:} The correspondence is functorial for extensions
of the base field.

\begin{center}\rule{0.5\linewidth}{0.5pt}\end{center}

\section{12. The Proof Strategy}\label{ch2002laurentlafforgue:the-proof-strategy}

\subsection{12.1 Step 1: Moduli Space of
Shtukas}\label{ch2002laurentlafforgue:step-1-moduli-space-of-shtukas}

Construct the moduli space \(\text{Cht}^r_k\) of shtukas of rank \(r\)
with \(k\) legs.

\subsection{12.2 Step 2: Cohomology}\label{ch2002laurentlafforgue:step-2-cohomology}

Compute the cohomology \(H^*(\text{Cht}^r_k)\) using the trace formula.
Show that the cohomology decomposes as a direct sum of automorphic
representations.

\subsection{12.3 Step 3: Hecke Operators}\label{ch2002laurentlafforgue:step-3-hecke-operators}

For each point \(x \in X\), the Hecke operator acts on the cohomology.
The eigenvalues of the Hecke operators at \(x\) are the Satake
parameters of the automorphic representation.

\subsection{12.4 Step 4: Galois
Representations}\label{ch2002laurentlafforgue:step-4-galois-representations}

The Frobenius morphism acts on the cohomology, giving an \(\ell\)-adic
representation of the Galois group of the function field. The
eigenvalues of Frobenius match the Hecke eigenvalues, establishing the
correspondence.

\subsection{12.5 Step 5:
Ramanujan--Petersson}\label{ch2002laurentlafforgue:step-5-ramanujanpetersson}

The Ramanujan--Petersson conjecture follows from Deligne's theorem on
the Weil conjectures.

\begin{center}\rule{0.5\linewidth}{0.5pt}\end{center}

\section{13. The Use of the Grothendieck--Lefschetz Trace
Formula}\label{ch2002laurentlafforgue:the-use-of-the-grothendiecklefschetz-trace-formula}

\subsection{13.1 The Formula}\label{ch2002laurentlafforgue:the-formula}

The Grothendieck--Lefschetz trace formula states:

\[
\sum_{x \in X(\mathbb{F}_{q^n})} \text{tr}(\text{Fr}^n, x) = \sum_{i=0}^{2d} (-1)^i \text{tr}(\text{Fr}^n, H^i_c(X))
\]

where the left side is the number of fixed points of the \(n\)-th power
of Frobenius, and the right side is the alternating sum of traces on
cohomology.

\subsection{13.2 Application to Shtukas}\label{ch2002laurentlafforgue:application-to-shtukas}

For the moduli space of shtukas, the left side counts shtukas over
\(\mathbb{F}_{q^n}\). Lafforgue shows that this count can be expressed
as:

\[
\sum_{\pi} \prod_{i=1}^k \text{tr}(\pi_x(\text{Fr}^n))
\]

where the sum is over automorphic representations \(\pi\) of GL(\(r\)).

\begin{center}\rule{0.5\linewidth}{0.5pt}\end{center}

\section{14. The Difficulty of the
Proof}\label{ch2002laurentlafforgue:the-difficulty-of-the-proof}

\subsection{14.1 Technical Challenges}\label{ch2002laurentlafforgue:technical-challenges}

The proof of the Langlands correspondence for GL(\(r\)) over function
fields involved overcoming numerous technical challenges:

\begin{enumerate}
\def\labelenumi{\arabic{enumi}.}
\item
 \textbf{Construction of the moduli space:} The moduli space of shtukas
 is a highly non-trivial algebraic stack. Its construction requires the
 theory of algebraic stacks and the theory of vector bundles on curves.
\item
 \textbf{Compactification:} The compactification of the moduli space
 required new techniques in toric geometry and the theory of algebraic
 stacks.
\item
 \textbf{The trace formula:} Applying the trace formula to the moduli
 space required a detailed analysis of the fixed points of Frobenius.
\item
 \textbf{The stabilization:} The stabilization of the trace formula is
 a major technical achievement in itself.
\item
 \textbf{The Hecke action:} Showing that the Hecke operators act on the
 cohomology in the expected way required deep results in the theory of
 perverse sheaves.
\end{enumerate}

\subsection{14.2 The Length of the Proof}\label{ch2002laurentlafforgue:the-length-of-the-proof}

Lafforgue's proof runs to several hundred pages, published in a series
of papers in \emph{Inventiones Mathematicae}. The main paper, ``Chtoucas
de Drinfeld et correspondance de Langlands,'' is over 300 pages long.

\begin{center}\rule{0.5\linewidth}{0.5pt}\end{center}

\section{15. The Function Field Case in
Context}\label{ch2002laurentlafforgue:the-function-field-case-in-context}

\subsection{15.1 Comparison with Number
Fields}\label{ch2002laurentlafforgue:comparison-with-number-fields}

The Langlands correspondence for GL(\(r\)) over number fields remains
one of the important open problems in mathematics. Lafforgue's
proof for function fields is a model for how the proof might proceed in
the number field case, but significant additional difficulties arise: -
\textbf{Archimedean places:} The behavior at infinite places requires
analysis. - \textbf{The trace formula:} The Arthur--Selberg trace
formula is more complicated for number fields. - \textbf{Modular forms:}
The theory of modular forms (for GL(2) over \(\mathbb{Q}\)) is the most
developed case.

\subsection{15.2 Recent Progress}\label{ch2002laurentlafforgue:recent-progress}

For number fields, the Langlands correspondence is known for: -
\textbf{GL(1):} Class field theory (proved by Takagi, Artin, etc.). -
\textbf{GL(2):} Proved by many authors (Wiles, Taylor, Kisin, et al.) in
the context of modularity. - \textbf{GL(n) for CM fields:} Proved by
Scholze, Shin, and others.

\begin{center}\rule{0.5\linewidth}{0.5pt}\end{center}

\section{16. The Impact of Lafforgue's
Work}\label{ch2002laurentlafforgue:the-impact-of-lafforgues-work}

\subsection{16.1 Impact on the Langlands
Program}\label{ch2002laurentlafforgue:impact-on-the-langlands-program}

Lafforgue's proof is the most significant result in the Langlands
program since Drinfeld's work on GL(2). It provides: - A complete proof
of the Langlands correspondence for GL(\(r\)) in the function field
case. - A model for the geometric approach to the Langlands program. -
Techniques that can be adapted to other groups.

\subsection{16.2 Impact on Algebraic
Geometry}\label{ch2002laurentlafforgue:impact-on-algebraic-geometry}

The construction of the moduli space of shtukas and its compactification
introduced new techniques in algebraic geometry: - The use of algebraic
stacks for moduli problems. - The compactification of moduli spaces
using toric methods. - The application of the trace formula to
cohomology.

\subsection{16.3 Impact on Number Theory}\label{ch2002laurentlafforgue:impact-on-number-theory}

The proof has implications for: - \textbf{The Ramanujan--Petersson
conjecture:} Proved in the function field case. - \textbf{The Sato--Tate
conjecture:} Follows from the Ramanujan--Petersson bounds. - \textbf{The
Langlands functoriality conjectures:} Provides evidence for the general
case.

\subsection{16.4 Impact on Representation
Theory}\label{ch2002laurentlafforgue:impact-on-representation-theory}

Lafforgue's work has influenced: - \textbf{The theory of automorphic
forms:} New techniques for studying automorphic representations. -
\textbf{The local Langlands correspondence:} The local correspondence
can be derived from the global one. - \textbf{The geometric Langlands
program:} The proof uses geometric methods that are central to the
geometric Langlands program.

\begin{center}\rule{0.5\linewidth}{0.5pt}\end{center}

\section{17. Lafforgue and Drinfeld}\label{ch2002laurentlafforgue:lafforgue-and-drinfeld}

\subsection{17.1 The Collaboration}\label{ch2002laurentlafforgue:the-collaboration}

Lafforgue's work builds directly on Drinfeld's earlier proof for GL(2).
Drinfeld introduced the concept of shtukas and proved the Langlands
correspondence for GL(2) over function fields.

\subsection{17.2 Drinfeld's Influence}\label{ch2002laurentlafforgue:drinfelds-influence}

Lafforgue acknowledges Drinfeld's influence: ``Without Drinfeld's work,
my own would have been impossible. His concept of shtukas is the
foundation on which everything is built.''

\begin{center}\rule{0.5\linewidth}{0.5pt}\end{center}

\section{18. The Work of Lafforgue's
Brother}\label{ch2002laurentlafforgue:the-work-of-lafforgues-brother}

\subsection{18.1 Vincent Lafforgue}\label{ch2002laurentlafforgue:vincent-lafforgue}

Laurent Lafforgue's younger brother, Vincent Lafforgue, is also a
distinguished mathematician. Vincent Lafforgue proved the Langlands
correspondence for GL(\(n\)) over function fields with additional
structures (the ``automorphic to Galois'' direction for all \(n\)).

\subsection{18.2 A Family of
Mathematicians}\label{ch2002laurentlafforgue:a-family-of-mathematicians}

The Lafforgue brothers are a rare example of two siblings both making
fundamental contributions to the same area of mathematics.

\begin{center}\rule{0.5\linewidth}{0.5pt}\end{center}

\section{19. Lafforgue's Other
Contributions}\label{ch2002laurentlafforgue:lafforgues-other-contributions}

\subsection{19.1 The Hitchin Fibration}\label{ch2002laurentlafforgue:the-hitchin-fibration}

Lafforgue has contributed to the theory of the Hitchin fibration, which
is a key tool in the geometric Langlands program. The Hitchin fibration
is a map from the moduli space of Higgs bundles to a vector space, and
its fibers are abelian varieties.

\subsection{19.2 The Theta
Correspondence}\label{ch2002laurentlafforgue:the-theta-correspondence}

Lafforgue has also worked on the theta correspondence, a relationship
between automorphic forms on dual reductive groups.

\begin{center}\rule{0.5\linewidth}{0.5pt}\end{center}

\section{20. The Proof of the Ramanujan--Petersson
Conjecture}\label{ch2002laurentlafforgue:the-proof-of-the-ramanujanpetersson-conjecture}

\subsection{20.1 Statement}\label{ch2002laurentlafforgue:statement}

The Ramanujan--Petersson conjecture for GL(\(r\)) over function fields
states that the Satake parameters of a cuspidal automorphic
representation have absolute value \(q^{1/2}\).

\subsection{20.2 Lafforgue's Proof}\label{ch2002laurentlafforgue:lafforgues-proof-1}

The proof follows from the Langlands correspondence: the Satake
parameters correspond to the eigenvalues of Frobenius on the associated
Galois representation, which are Weil numbers of weight \(r-1\) by
Deligne's theorem. This implies the Ramanujan--Petersson bounds.

\begin{center}\rule{0.5\linewidth}{0.5pt}\end{center}

\section{21. Later Work and Current
Directions}\label{ch2002laurentlafforgue:later-work-and-current-directions}

\subsection{21.1 Lafforgue's Work Since
2002}\label{ch2002laurentlafforgue:lafforgues-work-since-2002}

Since his Fields Medal, Lafforgue has continued to work on the Langlands
program and related areas: - \textbf{The geometric Langlands program:}
The categorical formulation and its connections to physics. -
\textbf{Shtukas for other groups:} Extending the theory to arbitrary
reductive groups. - \textbf{The Hitchin fibration and mirror symmetry.}

\subsection{21.2 The Global Langlands Correspondence for Arbitrary
Groups}\label{ch2002laurentlafforgue:the-global-langlands-correspondence-for-arbitrary-groups}

Lafforgue has been working on extending his results from GL(\(r\)) to
arbitrary reductive groups. This is the \textbf{Langlands functoriality}
problem, which remains one of the central open problems in the field.

\begin{center}\rule{0.5\linewidth}{0.5pt}\end{center}

\section{22. The Trace Formula and the
Stabilization}\label{ch2002laurentlafforgue:the-trace-formula-and-the-stabilization}

\subsection{22.1 The Arthur--Selberg Trace
Formula}\label{ch2002laurentlafforgue:the-arthurselberg-trace-formula}

The Arthur--Selberg trace formula is a non-abelian generalization of the
Poisson summation formula. For a reductive group \(G\) over a global
field \(F\), it states:

\[
\sum_{\pi} \text{tr}(\pi(f)) = \sum_{\gamma} J_\gamma(f)
\]

where the left side sums over automorphic representations \(\pi\) and
the right side sums over conjugacy classes \(\gamma\) in \(G(F)\).

\subsection{22.2 Lafforgue's
Stabilization}\label{ch2002laurentlafforgue:lafforgues-stabilization}

Lafforgue proved the \textbf{stabilization} of the trace formula for
GL(\(r\)) over function fields. This involves expressing the trace
formula as a sum over ``stable'' conjugacy classes, which are easier to
match with the geometry of the moduli space of shtukas.

\begin{center}\rule{0.5\linewidth}{0.5pt}\end{center}

\section{23. The Geometric Side of the Trace
Formula}\label{ch2002laurentlafforgue:the-geometric-side-of-the-trace-formula}

\subsection{23.1 Fixed Points of
Frobenius}\label{ch2002laurentlafforgue:fixed-points-of-frobenius}

For the moduli space of shtukas, the geometric side of the trace formula
counts the number of shtukas over finite fields. Lafforgue computed this
count explicitly:

\[
\#\text{Cht}^r_k(\mathbb{F}_{q^n}) = \sum_{\pi} \prod_{i=1}^k \text{tr}(\pi(\text{Fr}^n_{x_i}))
\]

\subsection{23.2 Matching with the Spectral
Side}\label{ch2002laurentlafforgue:matching-with-the-spectral-side}

The spectral side gives the sum over automorphic representations. By
equating the two sides, Lafforgue established the Langlands
correspondence.

\begin{center}\rule{0.5\linewidth}{0.5pt}\end{center}

\section{\texorpdfstring{24. The Rank \(r = 2\)
Case}{24. The Rank r = 2 Case}}\label{ch2002laurentlafforgue:the-rank-r-2-case}

\subsection{24.1 Drinfeld's Theorem}\label{ch2002laurentlafforgue:drinfelds-theorem}

For \(r = 2\), Drinfeld had already proved the Langlands correspondence.
His proof used: - \textbf{Drinfeld modules:} A generalization of
elliptic curves to function fields. - \textbf{Modular curves:} The
moduli space of Drinfeld modules. - \textbf{The Eichler--Shimura
congruence relation.}

\subsection{24.2 Lafforgue's
Generalization}\label{ch2002laurentlafforgue:lafforgues-generalization}

Lafforgue's proof for arbitrary \(r\) uses: - \textbf{Higher-dimensional
moduli spaces:} Not curves but algebraic stacks. - \textbf{Shtukas
instead of Drinfeld modules:} The correct generalization for higher
rank. - \textbf{The trace formula instead of the Eichler--Shimura
relation.}

\begin{center}\rule{0.5\linewidth}{0.5pt}\end{center}

\section{25. The Hitchin Fibration}\label{ch2002laurentlafforgue:the-hitchin-fibration-1}

\subsection{25.1 Definition}\label{ch2002laurentlafforgue:definition-1}

The \textbf{Hitchin fibration} is a map from the moduli space of Higgs
bundles to a vector space:

\[
h: \mathcal{M}_{Higgs}(r) \to \mathcal{A}_r
\]

where \(\mathcal{A}_r\) is the ``Hitchin base'' --- the space of
characteristic polynomials.

\subsection{25.2 Connection to Shtukas}\label{ch2002laurentlafforgue:connection-to-shtukas}

Lafforgue showed that the Hitchin fibration for GL(\(r\)) over a
function field is closely related to the moduli space of shtukas. This
connection provides an alternative approach to the Langlands
correspondence.

\begin{center}\rule{0.5\linewidth}{0.5pt}\end{center}

\section{26. The Global to Local
Correspondence}\label{ch2002laurentlafforgue:the-global-to-local-correspondence}

\subsection{26.1 The Local Langlands
Correspondence}\label{ch2002laurentlafforgue:the-local-langlands-correspondence}

The \textbf{local Langlands correspondence} relates: - Representations
of the Weil group of a local field. - Irreducible smooth representations
of GL(\(r\)) over the local field.

\subsection{26.2 Lafforgue's
Contribution}\label{ch2002laurentlafforgue:lafforgues-contribution-1}

The global Langlands correspondence (proved by Lafforgue) implies the
local correspondence for function fields. This is because the
restriction of a global Galois representation to a decomposition group
gives a local Galois representation.

\begin{center}\rule{0.5\linewidth}{0.5pt}\end{center}

\section{27. The Function Field Langlands Program: Other
Groups}\label{ch2002laurentlafforgue:the-function-field-langlands-program-other-groups}

\subsection{27.1 Beyond GL(r)}\label{ch2002laurentlafforgue:beyond-glr}

The Langlands program for function fields has been extended to other
groups: - \textbf{Vincent Lafforgue:} The Langlands correspondence for
GL(\(n\)) with arbitrary \(n\). - \textbf{Fleming and others:} Work on
symplectic and orthogonal groups.

\subsection{27.2 The Future}\label{ch2002laurentlafforgue:the-future}

The ultimate goal is to prove the Langlands correspondence for all
reductive groups over function fields. This would complete the function
field case of the Langlands program.

\begin{center}\rule{0.5\linewidth}{0.5pt}\end{center}

\section{28. The Work of Ngô Bảo
Châu}\label{ch2002laurentlafforgue:the-work-of-nguxf4-bux1ea3o-chuxe2u}

\subsection{28.1 The Fundamental Lemma}\label{ch2002laurentlafforgue:the-fundamental-lemma}

Ngô Bảo Châu proved the \textbf{fundamental lemma} for the trace
formula, a key ingredient in the stabilization of the Arthur--Selberg
trace formula. For this work, he was awarded the Fields Medal in 2010.

\subsection{28.2 Connection to Lafforgue}\label{ch2002laurentlafforgue:connection-to-lafforgue}

Lafforgue's proof of the Langlands correspondence does not require the
fundamental lemma (which was proved later), but the fundamental lemma is
essential for extending the proof to other groups and to the number
field case.

\begin{center}\rule{0.5\linewidth}{0.5pt}\end{center}

\section{29. Lafforgue's Impact on the Next
Generation}\label{ch2002laurentlafforgue:lafforgues-impact-on-the-next-generation}

\subsection{29.1 Training Students}\label{ch2002laurentlafforgue:training-students}

Lafforgue has supervised several PhD students and postdocs who have
continued to develop the theory of shtukas and the Langlands
correspondence.

\subsection{29.2 The IHÉS Seminar}\label{ch2002laurentlafforgue:the-ihuxe9s-seminar}

At IHÉS, Lafforgue has organized seminars and conferences on the
Langlands program, bringing together leading researchers in the field.

\begin{center}\rule{0.5\linewidth}{0.5pt}\end{center}

\section{30. Lafforgue's Philosophy of
Mathematics}\label{ch2002laurentlafforgue:lafforgues-philosophy-of-mathematics}

\subsection{30.1 Depth over Breadth}\label{ch2002laurentlafforgue:depth-over-breadth}

Lafforgue's work is characterized by its depth. Rather than working on
many problems, he focused on a single, significantly difficult problem
--- the Langlands correspondence --- and solved it completely for
GL(\(r\)) over function fields.

\subsection{30.2 The Unity of
Mathematics}\label{ch2002laurentlafforgue:the-unity-of-mathematics}

Lafforgue's proof exemplifies the unity of mathematics: it brings
together number theory (the Langlands program), algebraic geometry
(moduli spaces, stacks, cohomology), representation theory (automorphic
forms), and algebraic topology (the trace formula).

\begin{center}\rule{0.5\linewidth}{0.5pt}\end{center}

\section{31. The Deligne Conjectures and Lafforgue's
Work}\label{ch2002laurentlafforgue:the-deligne-conjectures-and-lafforgues-work}

\subsection{31.1 The Weil Conjectures}\label{ch2002laurentlafforgue:the-weil-conjectures}

The Weil conjectures, proved by Deligne in the 1970s, describe the
cohomology of algebraic varieties over finite fields. They state that
the eigenvalues of Frobenius are algebraic numbers of specific absolute
values (Weil numbers).

\subsection{31.2 Application to Shtukas}\label{ch2002laurentlafforgue:application-to-shtukas-1}

The moduli space of shtukas is an algebraic variety over a finite field,
so Deligne's theorem applies to its cohomology. The eigenvalues of
Frobenius on the cohomology are Weil numbers. By the trace formula,
these eigenvalues correspond to the Hecke eigenvalues of the associated
automorphic representations. This gives the Ramanujan--Petersson bounds.

\begin{center}\rule{0.5\linewidth}{0.5pt}\end{center}

\section{32. The Comparison with the Number Field
Case}\label{ch2002laurentlafforgue:the-comparison-with-the-number-field-case}

\subsection{32.1 Similarities}\label{ch2002laurentlafforgue:similarities}

The Langlands correspondence for number fields and function fields share
the same structure: - Galois representations correspond to automorphic
representations. - L-functions are preserved. - The local correspondence
is compatible.

\subsection{32.2 Differences}\label{ch2002laurentlafforgue:differences}

The proofs differ significantly: - \textbf{Function fields:} Geometric
methods (shtukas, cohomology) are available. - \textbf{Number fields:}
Analytic methods (trace formula, modular forms) are required. -
\textbf{Archimedean places:} Number fields have infinite places that
require analysis.

\subsection{32.3 The Future}\label{ch2002laurentlafforgue:the-future-1}

Lafforgue's proof serves as a model for the number field case. The key
challenge is to develop a geometric theory for number fields analogous
to the theory of shtukas.

\begin{center}\rule{0.5\linewidth}{0.5pt}\end{center}

\section{33. The Geometric Langlands
Program}\label{ch2002laurentlafforgue:the-geometric-langlands-program-1}

\subsection{33.1 The Categorical
Formulation}\label{ch2002laurentlafforgue:the-categorical-formulation}

The \textbf{geometric Langlands program}, developed by Beilinson and
Drinfeld, reformulates the Langlands correspondence as an equivalence of
categories:

\[
\text{D-modules on Bun}_G \leftrightarrow \text{Quasi-coherent sheaves on Loc}_{{}^L G}
\]

where: - Bun\(_G\) is the moduli stack of \(G\)-bundles on a curve. -
Loc\(_{{}^L G}\) is the moduli stack of local systems for the Langlands
dual group.

\subsection{33.2 Lafforgue's Connection}\label{ch2002laurentlafforgue:lafforgues-connection}

Lafforgue's work on shtukas is closely related to the geometric
Langlands program. The moduli space of shtukas is a geometric object
whose cohomology realizes the Langlands correspondence, connecting the
classical and geometric formulations.

\begin{center}\rule{0.5\linewidth}{0.5pt}\end{center}

\section{34. The Concept of Shtukas in
Detail}\label{ch2002laurentlafforgue:the-concept-of-shtukas-in-detail}

\subsection{34.1 Definition}\label{ch2002laurentlafforgue:definition-2}

A \textbf{shtuka} of rank \(r\) on a curve \(X\) over \(\mathbb{F}_q\)
is a tuple:

\[
\underline{\mathcal{E}} = (E, \sigma, \tau)
\]

where: - \(E\) is a vector bundle of rank \(r\) on \(X\). -
\(\sigma: E \to E\) is an isomorphism outside a finite set of points. -
\(\tau: \sigma^* E \to E\) is a morphism (the Frobenius structure).

\subsection{34.2 The Legs}\label{ch2002laurentlafforgue:the-legs}

The points where \(\sigma\) is not an isomorphism are called the
\textbf{legs} (or \textbf{poles}) of the shtuka. If there are \(k\)
legs, the shtuka is said to be a shtuka with \(k\) modifications.

\subsection{34.3 Moduli}\label{ch2002laurentlafforgue:moduli}

The moduli space of shtukas with \(k\) legs is an algebraic stack. Its
dimension is:

\[
\dim(\text{Cht}^r_k) = 2r(k-1)
\]

\begin{center}\rule{0.5\linewidth}{0.5pt}\end{center}

\section{35. The Cohomology of Shtuka
Spaces}\label{ch2002laurentlafforgue:the-cohomology-of-shtuka-spaces}

\subsection{35.1 The Betti Numbers}\label{ch2002laurentlafforgue:the-betti-numbers}

Lafforgue computed the Betti numbers of the moduli space of shtukas
using: - The trace formula. - The geometry of the moduli space. - The
theory of perverse sheaves.

\subsection{35.2 The Decomposition}\label{ch2002laurentlafforgue:the-decomposition}

The cohomology decomposes as:

\[
H^*(\text{Cht}^r_k) = \bigoplus_{\pi} \pi^{\otimes k}
\]

where the sum is over cuspidal automorphic representations \(\pi\) of
GL(\(r\)). This decomposition is the Langlands correspondence.

\begin{center}\rule{0.5\linewidth}{0.5pt}\end{center}

\section{36. Lafforgue's Mathematical
Style}\label{ch2002laurentlafforgue:lafforgues-mathematical-style}

\subsection{36.1 Systematic and Thorough}\label{ch2002laurentlafforgue:systematic-and-thorough}

Lafforgue's work is characterized by its systematic and thorough nature.
He leaves no gaps in his proofs and develops all the necessary machinery
from scratch.

\subsection{36.2 Geometric Intuition}\label{ch2002laurentlafforgue:geometric-intuition}

Despite the algebraic nature of the Langlands program, Lafforgue's work
is closely geometric. He thinks of automorphic forms in terms of the
geometry of moduli spaces, and his proofs are guided by geometric
intuition.

\subsection{36.3 Independence}\label{ch2002laurentlafforgue:independence}

Lafforgue worked largely independently on his proof. While building on
Drinfeld's ideas, he developed his own approach and solved problems that
had resisted solution by others.

\begin{center}\rule{0.5\linewidth}{0.5pt}\end{center}

\section{37. Impact on Arithmetic
Geometry}\label{ch2002laurentlafforgue:impact-on-arithmetic-geometry}

Lafforgue's work uses the Weil conjectures (proved by Deligne) in an
essential way. The eigenvalues of Frobenius on the cohomology of the
moduli space of shtukas are Weil numbers, which gives the
Ramanujan--Petersson bounds.

\subsection{22.2 The Grothendieck--Lefschetz
Formula}\label{ch2002laurentlafforgue:the-grothendiecklefschetz-formula}

The use of the Grothendieck--Lefschetz trace formula in the context of
shtukas is a model for how cohomological methods can be applied to
automorphic forms.

\begin{center}\rule{0.5\linewidth}{0.5pt}\end{center}

\section{Awards and Recognition}\label{ch2002laurentlafforgue:awards-and-recognition}

\textbf{2002:} - Fields Medal (International Mathematical Union)

\textbf{2001:} - Prix de la Société Mathématique de France

\textbf{2000:} - Prix Jacques Herbrand (Académie des Sciences)

\textbf{1999:} - Prix de l'Académie des Sciences

\begin{center}\rule{0.5\linewidth}{0.5pt}\end{center}

\section{Selected Major
Publications}\label{ch2002laurentlafforgue:selected-major-publications}

\begin{enumerate}
\def\labelenumi{\arabic{enumi}.}
\item
 \textbf{``Chtoucas de Drinfeld et correspondance de Langlands.''}
 \emph{Inventiones Mathematicae} 147, 1--241, 2002. --- \emph{The main
 paper proving the Langlands correspondence for GL(r) over function
 fields.}
\item
 \textbf{``Une compactification des champs de chtoucas.''}
 \emph{Journal of the American Mathematical Society} 15, 385--445,
 2002. --- \emph{The construction of the compactification of the moduli
 space of shtukas.}
\item
 \textbf{``Chtoucas de Drinfeld, formule des traces d'Arthur--Selberg
 et correspondance de Langlands.''} \emph{Proceedings of the
 International Congress of Mathematicians}, Beijing, 2002. --- \emph{An
 overview of the proof.}
\item
 \textbf{``La correspondance de Langlands pour les corps de
 fonctions.''} \emph{Séminaire Bourbaki}, 2003. --- \emph{An exposition
 of the Langlands correspondence for function fields.}
\item
 \textbf{``Pavages des espaces de chtoucas.''} \emph{Compositio
 Mathematica} 134, 1--57, 2002. --- \emph{The combinatorial structure
 of the moduli space of shtukas.}
\end{enumerate}

\begin{center}\rule{0.5\linewidth}{0.5pt}\end{center}

\begin{center}\rule{0.5\linewidth}{0.5pt}\end{center}

\section*{Exercises}
\addcontentsline{toc}{section}{Exercises}

\begin{exercise} Show that the number of points on an elliptic curve over $\\mathbb{F}_q$ can be expressed as $q + 1 - a_q$ where $a_q = \\alpha + \\bar{\\alpha}$ with $|\\alpha| = \\sqrt{q}$. \end{exercise}

\begin{exercise} Verify that the Langlands correspondence for $\\operatorname{GL}_2$ over $\\mathbb{F}_q(T)$ is a bijection between cuspidal automorphic representations and irreducible 2-dimensional Galois representations. \end{exercise}

\begin{exercise} Compute the L-function of a modular form of weight 2 and level 1. \end{exercise}

\begin{exercise} Prove that the Hecke operators $T_p$ commute and are simultaneously diagonalizable on the space of cusp forms. \end{exercise}

\begin{center}\rule{0.5\linewidth}{0.5pt}\end{center}

\section*{Summary}
\addcontentsline{toc}{section}{Summary}

Lafforgue proved the Langlands correspondence for function fields, establishing a deep bijection between automorphic representations of $\operatorname{GL}_n$ over function fields and Galois representations. His construction of the moduli space of shtukas and his use of étale cohomology provided the geometric framework for this fundamental result in number theory.

\section{References and Further
Reading}\label{ch2002laurentlafforgue:references-and-further-reading}

\subsection{Official Sources}\label{ch2002laurentlafforgue:official-sources}

\begin{itemize}
\tightlist
\item
 IMU Fields Medal 2002 Announcement:
 https://www.mathunion.org/imu-awards/fields-medal/fields-medal-2002
\item
 IHÉS: ``Laurent Lafforgue awarded Fields Medal.'' 2002.
\end{itemize}

\subsection{Key Technical References}\label{ch2002laurentlafforgue:key-technical-references}

\begin{itemize}
\tightlist
\item
 Lafforgue, L. ``Chtoucas de Drinfeld et correspondance de Langlands.''
 \emph{Inventiones Mathematicae} 147, 2002.
\item
 Drinfeld, V. ``Langlands conjecture for GL(2) over function fields.''
 \emph{Proceedings of the ICM}, 1978.
\item
 Laumon, G. ``La correspondance de Langlands sur les corps de
 fonctions.'' \emph{Séminaire Bourbaki}, 1996.
\item
 Laumon, G. and Moret-Bailly, L. \emph{Champs algébriques}. Springer,
 2000.
\end{itemize}

\subsection{Biography and Overview}\label{ch2002laurentlafforgue:biography-and-overview}

\begin{itemize}
\tightlist
\item
 Clozel, L. ``The Work of Laurent Lafforgue.'' ICM 2002 Laudatio.
\item
 Grothendieck, A. ``L'affaire Lafforgue.'' \emph{Gazette des
 Mathématiciens}, 2002.
\item
 Bernstein, J. and Gelbart, S. ``The Langlands program: an overview.''
 \emph{IAS Publication}, 2003.
\item
 Frenkel, E. ``The geometric Langlands program.'' \emph{Proceedings of
 the ICM}, 2010.
\end{itemize}

\subsection{Further Reading}\label{ch2002laurentlafforgue:further-reading}

\begin{itemize}
\tightlist
\item
 Deligne, P. \emph{Cohomologie étale}. Springer, 1977.
\item
 Laumon, G. and Moret-Bailly, L. \emph{Champs algébriques}. Springer,
 2000.
\item
 Drinfeld, V. ``Langlands conjecture for GL(2) over function fields.''
 \emph{Proceedings of the ICM}, 1978.
\item
 Arthur, J. ``The principle of functoriality.'' \emph{Bulletin of the
 AMS} 40, 2003.
\item
 Langlands, R. ``Problems in the theory of automorphic forms.''
 \emph{Lecture Notes in Mathematics} 170, Springer, 1970.
\end{itemize}

\begin{center}\rule{0.5\linewidth}{0.5pt}\end{center}

\emph{(End of report)}
